% GENERATED FILE. Edit canonical lessons, then run npm run generate:books.
% canonical-source-hash: fnv1a64:e626be89720cef09
% canonical-lessons: SA-C488-pratibudhyate, SA-C488-tandra, SA-C488-prabodhayati, SA-C488-unnidrah, SA-C488-praharah

\chapter{Sleep and Waking}
\label{ch:sa-sleep-and-waking}
\hlchaptermodality{488}

\begin{chapteropening}
\textbf{By the end of this chapter:} \emph{I can say how well someone slept.}

Complete the last lesson of chapter 488: I can say how well someone slept.
\end{chapteropening}

\section[\texorpdfstring{\sk{प्रतिबुध्यते} (pratibudhyate)}{प्रतिबुध्यते (pratibudhyate)}]{\sk{प्रतिबुध्यते} (pratibudhyate) — awakes}

\label{lesson:SA-C488-pratibudhyate}



\begin{quote}
Before the new one: say the Sanskrit for a scented cosmetic, then the Sanskrit for puts on clothes.
\end{quote}

\subsection*{You'll want to know: \sk{प्रतिबुध्यते}}
\textbf{\sk{प्रतिबुध्यते}} — \emph{pratibudhyate} — \textquotedblleft{}awakes\textquotedblright{}.

\begin{quote}
\sk{अरुणः} \sk{प्रातः} \sk{शीघ्रं} \sk{प्रतिबुध्यते।}
\end{quote}

\begin{quote}
\emph{aruṇaḥ prātaḥ śīghraṃ pratibudhyate.}
\end{quote}

\begin{quote}
Arun wakes early in the morning.
\end{quote}

Past (he …): \sk{सः} \sk{प्रातः} \sk{प्रतिबुद्धवान्}, \emph{saḥ prātaḥ pratibuddhavān}.

\begin{cousinweb}
\textbf{\sk{प्रतिबुध्यते}} (\emph{pratibudhyate}) is built from \emph{prati} and the root √budh, \textquotedblleft{}to wake, to know\textquotedblright{}, the root of \textbf{\sk{बुध्यते}}; MW gives \textquotedblleft{}to awaken, awake, wake\textquotedblright{} (MW).
\end{cousinweb}

\begin{grammarlens}[title={What you've built}]
One more word for describing what you lived and felt.
\end{grammarlens}

\subsection*{Your turn}
\begin{itemize}
\raggedright
  \item \emph{Say it:} \emph{pratibudhyate}
  \item \emph{Say it:} \emph{pratibudhyate}, once more
  \item \emph{Say it:} \emph{aruṇaḥ prātaḥ śīghraṃ pratibudhyate.} — Arun wakes early in the morning.
  \item \emph{Recall:} say the Sanskrit for a scented cosmetic, then the Sanskrit for puts on clothes, then say \emph{pratibudhyate} again
\end{itemize}

\begin{tcolorbox}[breakable,colback=teal!4,colframe=teal!35!black,title={Before you move on}]
What does \sk{प्रतिबुध्यते} mean? (\textquotedblleft{}Awakes\textquotedblright{}.) Say it once more.
\end{tcolorbox}
\section[\texorpdfstring{\sk{तन्द्रा} (tandrā)}{तन्द्रा (tandrā)}]{\sk{तन्द्रा} (tandrā) — lassitude; drowsiness}

\label{lesson:SA-C488-tandra}



\begin{quote}
Before the new one: say the Sanskrit for puts on clothes, then the Sanskrit for awakes.
\end{quote}

\subsection*{You'll want to know: \sk{तन्द्रा}}
\textbf{\sk{तन्द्रा}} — \emph{tandrā} — \textquotedblleft{}lassitude; drowsiness\textquotedblright{}.

\begin{quote}
\sk{भोजनात्} \sk{अनन्तरं} \sk{मम} \sk{तन्द्रा} \sk{भवति।}
\end{quote}

\begin{quote}
\emph{bhojanāt anantaraṃ mama tandrā bhavati.}
\end{quote}

\begin{quote}
After a meal I feel drowsy.
\end{quote}

A feminine noun.

\begin{cousinweb}
\textbf{\sk{तन्द्रा}} (\emph{tandrā}) MW gives \textquotedblleft{}lassitude, exhaustion, laziness\textquotedblright{}; \textquotedblleft{}drowsiness\textquotedblright{} is the sense it has in modern use; MW gives no derivation, so its origin is uncertain (MW).
\end{cousinweb}

\begin{grammarlens}[title={What you've built}]
One more word for describing what you lived and felt.
\end{grammarlens}

\subsection*{Your turn}
\begin{itemize}
\raggedright
  \item \emph{Say it:} \emph{tandrā}
  \item \emph{Say it:} \emph{tandrā}, once more
  \item \emph{Say it:} \emph{bhojanāt anantaraṃ mama tandrā bhavati.} — After a meal I feel drowsy.
  \item \emph{Recall:} say the Sanskrit for puts on clothes, then the Sanskrit for awakes, then say \emph{tandrā} again
\end{itemize}

\begin{tcolorbox}[breakable,colback=teal!4,colframe=teal!35!black,title={Before you move on}]
What does \sk{तन्द्रा} mean? (\textquotedblleft{}Lassitude; drowsiness\textquotedblright{}.) Say it once more.
\end{tcolorbox}
\section[\texorpdfstring{\sk{प्रबोधयति} (prabodhayati)}{प्रबोधयति (prabodhayati)}]{\sk{प्रबोधयति} (prabodhayati) — wakes someone up}

\label{lesson:SA-C488-prabodhayati}



\begin{quote}
Before the new one: say the Sanskrit for awakes, then the Sanskrit for lassitude.
\end{quote}

\subsection*{You'll want to know: \sk{प्रबोधयति}}
\textbf{\sk{प्रबोधयति}} — \emph{prabodhayati} — \textquotedblleft{}wakes someone up\textquotedblright{}.

\begin{quote}
\sk{माता} \sk{प्रातः} \sk{पुत्रं} \sk{प्रबोधयति।}
\end{quote}

\begin{quote}
\emph{mātā prātaḥ putraṃ prabodhayati.}
\end{quote}

\begin{quote}
Mother wakes her son in the morning.
\end{quote}

Past (she …): \sk{सा} \sk{पुत्रं} \sk{प्रबोधितवती}, \emph{sā putraṃ prabodhitavatī}.

\begin{cousinweb}
\textbf{\sk{प्रबोधयति}} (\emph{prabodhayati}) is built on \emph{pra} and the causative of the root √budh, \textquotedblleft{}to wake, to know\textquotedblright{}, the root of \textbf{\sk{बुध्यते}}: to make someone wake; MW gives \textquotedblleft{}to wake up, awaken (trans.)\textquotedblright{} (MW).
\end{cousinweb}

\begin{grammarlens}[title={What you've built}]
One more word for describing what you lived and felt.
\end{grammarlens}

\subsection*{Your turn}
\begin{itemize}
\raggedright
  \item \emph{Say it:} \emph{prabodhayati}
  \item \emph{Say it:} \emph{prabodhayati}, once more
  \item \emph{Say it:} \emph{mātā prātaḥ putraṃ prabodhayati.} — Mother wakes her son in the morning.
  \item \emph{Recall:} say the Sanskrit for awakes, then the Sanskrit for lassitude, then say \emph{prabodhayati} again
\end{itemize}

\begin{tcolorbox}[breakable,colback=teal!4,colframe=teal!35!black,title={Before you move on}]
What does \sk{प्रबोधयति} mean? (\textquotedblleft{}Wakes someone up\textquotedblright{}.) Say it once more.
\end{tcolorbox}
\section[\texorpdfstring{\sk{उन्निद्रः} (unnidraḥ)}{उन्निद्रः (unnidraḥ)}]{\sk{उन्निद्रः} (unnidraḥ) — sleepless}

\label{lesson:SA-C488-unnidrah}



\begin{quote}
Before the new one: say the Sanskrit for lassitude, then the Sanskrit for wakes someone up.
\end{quote}

\subsection*{You'll want to know: \sk{उन्निद्रः}}
\textbf{\sk{उन्निद्रः}} — \emph{unnidraḥ} — \textquotedblleft{}sleepless\textquotedblright{}.

\begin{quote}
\sk{उद्वेगेन} \sk{पिता} \sk{रात्रौ} \sk{उन्निद्रः} \sk{भवति।}
\end{quote}

\begin{quote}
\emph{udvegena pitā rātrau unnidraḥ bhavati.}
\end{quote}

\begin{quote}
Father lies sleepless at night with worry.
\end{quote}

\begin{cousinweb}
\textbf{\sk{उन्निद्रः}} (\emph{unnidraḥ}) is built from \emph{ud}, \textquotedblleft{}up, out\textquotedblright{}, and \textbf{\sk{निद्रा}}, \textquotedblleft{}sleep\textquotedblright{}; MW gives \textquotedblleft{}sleepless, awake\textquotedblright{} (MW).
\end{cousinweb}

\begin{grammarlens}[title={What you've built}]
One more word for describing what you lived and felt.
\end{grammarlens}

\subsection*{Your turn}
\begin{itemize}
\raggedright
  \item \emph{Say it:} \emph{unnidraḥ}
  \item \emph{Say it:} \emph{unnidraḥ}, once more
  \item \emph{Say it:} \emph{udvegena pitā rātrau unnidraḥ bhavati.} — Father lies sleepless at night with worry.
  \item \emph{Recall:} say the Sanskrit for lassitude, then the Sanskrit for wakes someone up, then say \emph{unnidraḥ} again
\end{itemize}

\begin{tcolorbox}[breakable,colback=teal!4,colframe=teal!35!black,title={Before you move on}]
What does \sk{उन्निद्रः} mean? (\textquotedblleft{}Sleepless\textquotedblright{}.) Say it once more.
\end{tcolorbox}
\section[\texorpdfstring{\sk{प्रहरः} (praharaḥ)}{प्रहरः (praharaḥ)}]{\sk{प्रहरः} (praharaḥ) — a three-hour watch}

\label{lesson:SA-C488-praharah}



\begin{quote}
Before the new one: say the Sanskrit for wakes someone up, then the Sanskrit for sleepless.
\end{quote}

\subsection*{You'll want to know: \sk{प्रहरः}}
\textbf{\sk{प्रहरः}} — \emph{praharaḥ} — \textquotedblleft{}a three-hour watch\textquotedblright{}.

\begin{quote}
\sk{रात्रेः} \sk{प्रथमे} \sk{प्रहरे} \sk{सः} \sk{स्वपिति।}
\end{quote}

\begin{quote}
\emph{rātreḥ prathame prahare saḥ svapiti.}
\end{quote}

\begin{quote}
He sleeps in the first watch of the night.
\end{quote}

A masculine noun.

\begin{cousinweb}
\textbf{\sk{प्रहरः}} (\emph{praharaḥ}) is built from \emph{pra} and the root √hṛ, \textquotedblleft{}to take, to strike\textquotedblright{}; MW gives \textquotedblleft{}the eighth part of a day, a watch\textquotedblright{} (MW).
\end{cousinweb}

\begin{grammarlens}[title={What you've built}]
One more word for describing what you lived and felt.
\end{grammarlens}

\subsection*{Your turn}
\begin{itemize}
\raggedright
  \item \emph{Say it:} \emph{praharaḥ}
  \item \emph{Say it:} \emph{praharaḥ}, once more
  \item \emph{Say it:} \emph{rātreḥ prathame prahare saḥ svapiti.} — He sleeps in the first watch of the night.
  \item \emph{Recall:} say the Sanskrit for wakes someone up, then the Sanskrit for sleepless, then say \emph{praharaḥ} again
\end{itemize}

\begin{tcolorbox}[breakable,colback=teal!4,colframe=teal!35!black,title={Before you move on}]
What does \sk{प्रहरः} mean? (\textquotedblleft{}A three-hour watch\textquotedblright{}.) Say it once more.
\end{tcolorbox}
